\chapterheading{5. প্রোগ্রামিং ভাষা}

\begin{overviewbox}

কম্পিউটার নিজে থেকে কোনো সমস্যা সমাধান করতে পারে না।
কোনো নির্দিষ্ট কাজ সম্পন্ন করতে তাকে ধাপে ধাপে নির্দেশ দিতে হয়।
এই নির্দেশনার সমষ্টিকে \textbf{Program} এবং Program লেখার জন্য
ব্যবহৃত নিয়মবদ্ধ ভাষাকে \textbf{Programming Language} বলা হয়।

Programming Language-এর বিকাশ হয়েছে কম্পিউটারের
Machine Language থেকে মানুষের জন্য তুলনামূলক সহজ
High Level Language পর্যন্ত। এই অধ্যায়ে Machine Language,
Assembly Language, Middle Level Language, High Level Language
এবং Fourth Generation Language-এর ধারণা ও বৈশিষ্ট্য জানা যাবে।
একই সঙ্গে Compiler, Assembler ও Interpreter-এর মতো
\textbf{Translator Program} কীভাবে Source Program-কে
কম্পিউটারের উপযোগী রূপে রূপান্তর করে, তাও শেখা হবে।

Program তৈরি শুধু Code লেখার কাজ নয়। একটি ভালো Program তৈরির আগে
Problem বুঝতে হয়, প্রয়োজনীয় Input ও Output নির্ধারণ করতে হয় এবং
সমাধানের Process ও Logic পরিকল্পনা করতে হয়। তাই এ অধ্যায়ে
\textbf{Program Organization, Program Design Model, Modular Programming,
Computational Thinking, Problem Analysis ও IPO Model} আলোচনা করা হবে।

এরপর Program Development-এর বিভিন্ন ধাপ, Dry Run, Trace Table,
Algorithm এবং Flowchart-এর মাধ্যমে Program Logic তৈরি, প্রকাশ ও
পরীক্ষা করার পদ্ধতি শেখানো হবে।

অধ্যায়ের পরবর্তী অংশে থাকবে \textbf{Python Programming Language}।
Python-এর পরিচিতি, বৈশিষ্ট্য, Programming Environment, Program লেখা
ও চালানো, Syntax, Variable, Data Type, Input--Output এবং Operator
সম্পর্কে শেখা হবে। এরপর Conditional Logic,
\textenglish{\ttfamily for} Loop,
\textenglish{\ttfamily while} Loop,
Loop Control Statement এবং Loop Pattern ব্যবহার করে বিভিন্ন
Problem সমাধান করা হবে।

শেষ অংশে Data Structure ও Function-এর মৌলিক ধারণা থাকবে।
Function ব্যবহারের পাশাপাশি \textbf{Function Contract, Function Scope,
Local ও Global Variable} এবং Program-এ \textbf{File কেন প্রয়োজন ও Python-এ
File Handling কীভাবে করা হয়} সে বিষয়েও ধারণা পাওয়া যাবে।

এই অধ্যায়ের মূল লক্ষ্য শুধু Programming Syntax শেখা নয়;
বরং একটি Problem বুঝে, বিশ্লেষণ করে, Logic তৈরি করে এবং
সঠিকভাবে Program-এর মাধ্যমে সমাধান করতে শেখা।

\end{overviewbox}

\begin{learningbox}

\begin{itemize}
\item Program ও Programming Language-এর ধারণা ব্যাখ্যা করতে।
\item Machine Language, Assembly Language, Middle Level Language, High Level Language ও 4GL-এর বৈশিষ্ট্য বুঝতে।
\item Compiler, Assembler ও Interpreter-এর কাজ ও পার্থক্য ব্যাখ্যা করতে।
\item Program Organization, Program Design Model ও Modular Programming-এর ধারণা বুঝতে।
\item Computational Thinking ব্যবহার করে একটি Problem বিশ্লেষণ করতে।
\item Input, Process ও Output শনাক্ত করে IPO Model তৈরি করতে।
\item Program Development-এর বিভিন্ন ধাপ ব্যাখ্যা করতে।
\item Dry Run, Trace Table, Algorithm ও Flowchart ব্যবহার করে Program Logic যাচাই ও প্রকাশ করতে।
\item Python-এর মৌলিক Syntax, Variable, Data Type, Input--Output ও Operator ব্যবহার করতে।
\item Conditional Logic, \textenglish{\ttfamily for} Loop এবং \textenglish{\ttfamily while} Loop ব্যবহার করে Problem সমাধান করতে।
\item Loop Control Statement ও Loop Pattern প্রয়োগ করতে।
\item Data Structure ও Function-এর মৌলিক ধারণা বুঝতে।
\item Function Contract, Function Scope এবং Local ও Global Variable-এর পার্থক্য ব্যাখ্যা করতে।
\item Program-এ File ব্যবহারের প্রয়োজনীয়তা বুঝতে এবং Python-এ File Open, Read, Write ও Append-এর মৌলিক ব্যবহার ব্যাখ্যা করতে।
\item Python ব্যবহার করে সহজ বাস্তব Problem-এর Program তৈরি করতে।
\end{itemize}

\end{learningbox}

\begin{pathbox}
\textbf{ধাপ ১: প্রোগ্রামিংয়ের ভিত্তি}\\
Program $\rightarrow$ Programming Language $\rightarrow$ Language Level $\rightarrow$ Translator Program $\rightarrow$ Program Design

\vspace{0.5mm}
$\Downarrow$

\textbf{ধাপ ২: সমস্যা বিশ্লেষণ ও সমাধানের পরিকল্পনা}\\
Computational Thinking $\rightarrow$ Problem Analysis $\rightarrow$ IPO Model $\rightarrow$ Program Development $\rightarrow$ Algorithm \& Flowchart

\vspace{0.5mm}
$\Downarrow$

\textbf{ধাপ ৩: Python প্রোগ্রামিং শেখা}\\
Python $\rightarrow$ Variable \& Data Type $\rightarrow$ Input/Output $\rightarrow$ Operator $\rightarrow$ Conditional Logic

\vspace{0.5mm}
$\Downarrow$

\textbf{ধাপ ৪: উন্নত প্রোগ্রামিং ধারণা}\\
Loop $\rightarrow$ Data Structure $\rightarrow$ Function $\rightarrow$ File Handling
\end{pathbox}

\small

%================================================
% 5.1
%================================================
